\documentclass[../../main.tex]{subfiles}
\begin{document}


{\bf [解]}

数列
\begin{align}
 \begin{cases}
  a_n&=\displaystyle\int_0^1t^{2n-1}e^t\,dt \\
  b_n&=\displaystyle\sum_{k=1}^{n-1}\dfrac{{}_{2n-1}P_{2n-2k}}{2k+1}
 \end{cases} \label{eq:1}
\end{align}
とおくと題意の関係式は
\begin{align}
a_n+e\,b_n=(2n-1)! \quad\cdots\Diamond \label{eq:2}
\end{align}
である．以下$n\ge 2$に対して\cref{eq:2}を数学的帰納法で示す．

$n=2$の時は\cref{eq:1}より
\begin{align}
\begin{cases}
a_2&=\int_0^1t^3e^t\,dt=\left[e^t(t^3-3t^2+6t-6)\right]_0^1=(-2e)-(-6)=6-2e \\
b_2&=\dfrac{1}{3}\,{}_3P_2=\dfrac{1}{3}\cdot3\cdot2=2 
\end{cases}
\end{align}
だから$a_2+eb_2=6-2e+2e=6=3!$となり\cref{eq:2}は$n=2$で成立する．

そこで以下$n=m\ (m\ge2)$での\cref{eq:2}の成立を仮定して$n=m+1$での成立を示す．部分積分を2回用いて$a_{m+1}$を変形すると
\begin{align}
a_{m+1}
 &=\int_0^1t^{2m+1}e^t\,dt \\
 &=\left[t^{2m+1}e^t\right]_0^1-(2m+1)\int_0^1t^{2m}e^t\,dt \\
 &=e-(2m+1)\left\{\left[t^{2m}e^t\right]_0^1-2m\int_0^1t^{2m-1}e^t\,dt\right\} \\
 &=e-(2m+1)\{e-2m\,a_m\} \\
 \therefore\ 
a_{m+1}&=-2m\,e+(2m+1)(2m)a_m \quad\cdots\text{①} \label{eq:3}
\end{align}
である．

次に$b_{m+1}$について，順列の定義${}_nP_r=n!/(n-r)!$から${}_{2m-1}P_{2m-2k}=(2m-1)!/(2k-1)!$であり，
\begin{align}
{}_{2m+1}P_{2m+2-2k}=\dfrac{(2m+1)!}{(2k-1)!}=(2m+1)(2m)\cdot\dfrac{(2m-1)!}{(2k-1)!}=(2m+1)(2m)\cdot{}_{2m-1}P_{2m-2k}
\end{align}
だから
\begin{align}
b_{m+1}=\sum_{k=1}^m\dfrac{{}_{2m+1}P_{2m+2-2k}}{2k+1}=(2m+1)(2m)\sum_{k=1}^m\dfrac{{}_{2m-1}P_{2m-2k}}{2k+1}
\end{align}
となる．
和を$k=m$の項（${}_{2m-1}P_0/(2m+1)=1/(2m+1)$）と残り（$=b_m$）に分けて，
\begin{align}
b_{m+1}=(2m+1)(2m)\left\{b_m+\dfrac{1}{2m+1}\right\}=2m+(2m+1)(2m)b_m \quad\cdots\text{②} \label{eq:4}
\end{align}
とかける．

\cref{eq:3,eq:4}から
\begin{align}
\begin{split}
 a_{m+1}+e\,b_{m+1}
  &=\{-2me+(2m+1)(2m)a_m\}+e\{2m+(2m+1)(2m)b_m\} \\
  &=(2m+1)(2m)(a_m+e\,b_m)=(2m+1)(2m)(2m-1)!\quad(\because\text{仮定})=(2m+1)!
\end{split}
\end{align}
となり，$n=m+1$でも\cref{eq:2}が成立する．

以上から数学的帰納法により\cref{eq:2}は全ての$n\ge2$で成立する．


{\bf [解2]}
解答では帰納法を利用したが，もっと直接的に数列の一般項を求めることもできるのでその方法も紹介する．
というのも，実は\cref{eq:3}から$a_n$の一般項を求められるからだ．\cref{eq:3}の両辺を$(2m+1)!$で割って$m$を$n$で書き直すと
\begin{align}
\dfrac{a_{n+1}}{(2n+1)!}=\dfrac{a_n}{(2n-1)!}-\dfrac{2n}{(2n+1)!}\,e=\dfrac{a_n}{(2n-1)!}-e\left\{\dfrac{1}{(2n)!}-\dfrac{1}{(2n+1)!}\right\}
\end{align}
だからこれを$n=1,2,\cdots$と繰り返し用いて
\begin{align}
 \dfrac{a_n}{(2n-1)!}=a_1-e\sum_{k=1}^{n-1}\left\{\dfrac{1}{(2k)!}-\dfrac{1}{(2k+1)!}\right\} 
\end{align}
である．
ここで$a_1$は
\begin{align}
 a_1=\int_0^1te^t\,dt=\left[e^t(t-1)\right]_0^1=0-(-1)=1
\end{align}
だから代入して
\begin{align}
 \dfrac{a_n}{(2n-1)!}=1-e\sum_{k=1}^{n-1}\left\{\dfrac{1}{(2k)!}-\dfrac{1}{(2k+1)!}\right\} \\
\therefore
  a_n + e(2n-1)!\sum_{k=1}^{n-1}\left\{\dfrac{1}{(2k)!}-\dfrac{1}{(2k+1)!}\right\} = (2n-1)! \label{eq:5} 
\end{align}
である．
\cref{eq:5}と\cref{eq:2}を比較すると
\begin{align}
 b_n = e(2n-1)!\sum_{k=1}^{n-1}\left\{\dfrac{1}{(2k)!}-\dfrac{1}{(2k+1)!}\right\} \label{eq:6}
\end{align}
であれば題意は示されるから以下これを示す．

${}_{2n-1}P_{2n-2k}=(2n-1)!/(2k-1)!$だから
\begin{align}
b_n
 &=(2n-1)!\sum_{k=1}^{n-1}\dfrac{1}{(2k+1)(2k-1)!} \\
 &=(2n-1)!\sum_{k=1}^{n-1}\dfrac{2k}{(2k+1)!} \\
 &=(2n-1)!\sum_{k=1}^{n-1}\left\{\dfrac{1}{(2k)!}-\dfrac{1}{(2k+1)!}\right\} \label{eq:7}
\end{align}
となり，\cref{eq:6}と一致する．よって示された．


\newpage

\end{document}
